%%%%%%%%%%%%%%%%%%%%%%%%%%%%%%%%%%%%%%%%%%%%%%%%%%%%%%%%%%%%%%%%%%%%%%%%%%%
%%  Project: CPHOS LaTeX Template
%%  File: example-paper.tex
%%%%%%%%%%%%%%%%%%%%%%%%%%%%%%%%%%%%%%%%%%%%%%%%%%%%%%%%%%%%%%%%%%%%%%%%%%%

\documentclass[exam]{../cphos-e}

\cphostitle{第XX届 CPHOS 物理竞赛联考}
\cphossubtitle{实验试题}
\cphoswatermark{../assets/watermark.jpg}

\begin{document}

\maketitlepage

{\kaiti\small 本试题于 2026 年 3 月 9 日发布，最后更新于\compiletime 。}

{\kaiti\small CPHOS 物理竞赛联考是开放性公益性的考试，有意向参与的教师和学生可以关注“CPHOS”微信公众号报名。}

\begin{examschedule}
\begin{tabular}{@{}cc@{}}
    答题卡上传 & 阅卷 \\
    YYYY/MM/DD HH:MM - YYYY/MM/DD HH:MM & YYYY/MM/DD HH:MM - YYYY/MM/DD HH:MM \\
\end{tabular}

\vspace{0.3em}

\begin{tabular}{@{}ccc@{}}
    非正式成绩 & 成绩申诉 & 正式成绩 \\
    YYYY/MM/DD HH:MM & YYYY/MM/DD HH:MM - YYYY/MM/DD HH:MM & YYYY/MM/DD HH:MM \\
\end{tabular}
\end{examschedule}

\begin{examnotes}
    \item 实验试题共 \textbf{\underline{\cphostotalpages}} 页，实验答题卡共 \textbf{\underline{3}} 页，答题时间 \textbf{\underline{60}} 分钟，试题满分 \textbf{\underline{\cphostotalscore}} 分。
    \item 请在答题卡指定区域作答，草稿纸内容不计分。
    \item 若发现试题表述问题，请向领队（教练）反馈。
\end{examnotes}

\vspace{1em}

\begin{problem}[40]{弗兰克-赫兹实验}
\begin{expcontent}
1913 年玻尔模型指出原子能级是分立的。弗兰克-赫兹实验通过电子与气体原子碰撞电离/激发行为，验证了原子能级量子化。

\begin{figure}[H]
    \centering
    \fbox{\parbox[c][5.5cm][c]{0.72\textwidth}{\centering 实验装置图占位示意}}
    \caption{弗兰克-赫兹实验装置示意图}
    \label{fig:fh-device}
\end{figure}

\pmark{观察与定性分析}
\subq{在固定灯丝加热电压与栅极电压条件下，调节加速电压并记录电流变化特征。}
\subsubq{说明首次电流峰值出现前后电子与汞原子碰撞行为的变化。}
\subsubq{讨论异常下降段可能对应的实验误差来源，并给出修正建议。}

\pmark{数据处理与参数估计}
\subq{根据实验表格中相邻峰值间隔，估算第一激发能对应电势差。}
\subsubq{结合误差传播，给出测量结果的不确定度表达。}
\subsubq{比较实验值与参考值，说明是否在误差范围内一致。}
\end{expcontent}
\end{problem}

\begin{copyrightinfo}
\begin{center}
\small
Copyright \copyright\ \the\year\ CPHOS. All Rights Reserved.
\end{center}
\end{copyrightinfo}

\end{document}
